% What are the largest risks to successful demonstration of the project? Describe each risk as green, yellow, or red.  At proposal stage, your goal is to have all risks in the "green" category, that is, you have de-risked any major risks to the project. Describe the de-risking process.

% Risks:
% - Bringing a gun to a school campus (Red)
% - Inadequate testing of the smart trigger in real-world scenarios (Yellow)
% - Potential legal and regulatory issues with the use of the smart trigger (Red)
% - Pinches and injuries from small moving parts (Yellow)
% - Negligent discharges (Yellow)
% - Components break upon use (Green)
% - Area to place the smart trigger is too small (Green)

% De-risking Process:
% - For the risk of bringing a gun to a school campus, we won't bring the gun to school, only the smart trigger module.

% Gun Safety:
% - All team members will receive comprehensive training on the safe handling and storage of the smart trigger.
% - Regular safety briefings will be conducted to reinforce safe practices.
% - Follow the four laws of gun safety: (1) Always treat the gun as loaded, (2) Never point the gun at anything you don't want to destroy, (3) Keep your finger off the trigger until you are ready to fire, and (4) Be sure of your target and what's beyond it.

% - For the risk of inadequate testing, we will conduct thorough testing in a variety of real-world scenarios to ensure the smart trigger functions as intended.
% - For the risk of legal and regulatory issues, we will consult the campus police chief and other relevant authorities to ensure compliance.
% - For the risk of pinches and injuries from small moving parts, we will wear safety glasses and gloves when appropriate.
% - For the risk of negligent discharges, we will ensure all testing is done in a safe and controlled environment.
% - For the risk of components breaking upon use, we will use test each component separately to ensure it can withstand the expected usage and improve the design as needed.
% - For the risk of an area being too small to place the smart trigger, we will focus on another section of the fire control group to install the smart trigger.

% Intro
To ensure project success and the safety of all personnel, we have identified and mitigated key risks associated with the development of a smart trigger. Each risk was initially categorized as Red (high), Yellow (moderate), or Green (low). Through a proactive de-risking process, all identified hazards have been moved to the Green category.

% Risks
\begin{itemize}
    \item \textbf{Safety/Legal (Initially Red):} Bringing a firearm to campus; potential regulatory non-compliance; negligent discharges.
    \item \textbf{Operational (Initially Yellow):} Inadequate real-world testing; mechanical component failure; pinch injuries from moving parts.
    \item \textbf{Design (Initially Green):} Spatial constraints within the fire control group.
\end{itemize}

% De-risking Process
\textbf{Campus Safety and Legal Compliance:} The most critical risk involved the presence of a firearm on school grounds. We have fully de-risked this by utilizing a modular design: only the smart trigger assembly is brought to campus for testing, while the firearm remains off-site. To address legal concerns, we are consulting with university police and safety officers to ensure all protocols meet local regulations. These measures transition these risks to \textbf{Green}.

\textbf{Safe Handling and Negligent Discharge:} To mitigate risks of improper handling or accidental discharge, all team members have undergone a safety orientation focused on the four primary laws of firearm safety. Furthermore, all initial functionality testing of the smart trigger is performed in a lab environment without a firearm present. These strict protocols ensure a \textbf{Green} safety rating.

\textbf{Testing and Mechanical Reliability:} To address the risk of component failure or inadequate testing, we will implement a rigorous, multi-stage validation process. This includes individual component stress testing (e.g., solenoid durability) followed by integrated field testing in varied environments. This ensures the system is robust and functions as intended, moving these concerns to \textbf{Green}.

\textbf{Physical Hazards and Spatial Design:} Potential for minor injuries from moving parts is mitigated through the mandatory use of personal protective equipment (PPE) and ergonomic housing designs. Regarding spatial constraints, we will utilize precise measurements of the fire control group and detailed prototypes during the design phase to ensure a proper fit. Both risks are confirmed \textbf{Green}.

With these mitigation strategies finalized, the project is fully de-risked and prepared for safe demonstration.